\chapter{Simulation and Complete Search}\label{ch:search}

The simplest correct algorithm is to try everything. When the input is small enough, that is also the intended solution;
when it is not, a complete search is still where you begin: it pins down exactly what you are computing, it is the
reference you stress test against, and it often earns the points of the small subtasks. This chapter is about doing
exhaustive work \emph{systematically}: simulating a process exactly as described, enumerating subsets and permutations
without missing or repeating any, and organising a search so that it abandons hopeless branches early. It closes with three
techniques that stretch the reach of exhaustive search far beyond its naive limits: meet in the middle, fracturing search,
and the brute-force-assisted casework that solves so many ad hoc problems.

\begin{chapterintro}{\bronze\ \silver\ \gold\ \ioi}{Chapters~\ref{ch:approach} to~\ref{ch:cpp}.}
  \item simulate processes correctly, and skip ahead when a state repeats;
  \item estimate the size of a search space and choose the smallest one;
  \item enumerate subsets (bitmasks, Gray codes, submasks, $k$-subsets) and permutations;
  \item write backtracking searches and prune them by orders of magnitude;
  \item halve an exponent with meet in the middle, and list the $k$ best solutions with fracturing search;
  \item solve casework problems by tabulating small cases with a brute force.
\end{chapterintro}

% ======================================================================================================================
\section[Simulation]{Simulation \hfill \bronze}\label{sec:simulation}

Some problems ask you to do exactly what the statement describes: move a robot for $T$ steps, apply a list of swaps, play a
game by fixed rules, run a queue of customers through a shop. There is no clever algorithm to discover; the whole difficulty
is in getting every detail right, and a handful of habits turn that from a gamble into a routine.

\begin{itemize}
  \item \textbf{Make the state explicit.} Write down exactly what describes the situation at one moment (positions,
        directions, counters, whose turn it is) and keep it in a few named variables or a struct. Everything the rules
        mention must be part of the state, and nothing else should be.
  \item \textbf{Write one step function.} A single piece of code that advances the state by one step is far easier to
        check than logic spread across a loop. Test it on the first steps of the sample before running it for long.
  \item \textbf{Use direction arrays for grids.} Store the four directions in clockwise order (up, right, down, left) as
        \cc|dr = {-1, 0, 1, 0}| and \cc|dc = {0, 1, 0, -1}|. Then moving forward is \cc{r += dr[d], c += dc[d]}, turning
        right is \cc{d = (d + 1) % 4}, turning left \cc{d = (d + 3) % 4} and turning around \cc{d ^= 2}. Eight directions
        work the same way with arrays of length eight.
  \item \textbf{Pad the grid.} Surrounding the grid with a border of blocked cells removes every bounds check: a step off
        the grid simply hits a wall. The Grid Paths solution in Section~\ref{sec:backtracking} uses this.
  \item \textbf{Trace the samples.} Print the state after every step with \cc{dbg} and compare it with a hand simulation of
        the first sample; the first line that differs is where the bug is.
\end{itemize}

\subsection{When $T$ is huge: find the cycle}

If the number of possible states is finite, a deterministic process must eventually revisit a state, and from then on it
repeats forever, because the next state depends only on the current one. So a process that runs for $T = 10^{18}$ steps is
never simulated step by step. Record the time at which each state first appears; when a state appears again at time $t$,
having first appeared at time $s$, the process has a \defn[cycle detection]{cycle} of length $L = t - s$, and the state at
any time $T \ge s$ equals the state at time $s + (T - s) \bmod L$. Equivalently, jump over $\lfloor (T - t) / L\rfloor$
whole cycles at once and simulate the fewer than $L$ steps that remain. Drawn as a sequence of states, the process has the
shape of the Greek letter $\rho$: a \emph{tail} of states that are never seen again, leading into the cycle
(Figure~\ref{fig:robot}).

\begin{problem}{Robot on a grid}
A robot stands on a grid with walls, facing up. At each step it moves forward if the cell in front is free and inside the
grid; otherwise it turns right (and does not move in that step). Where is it after $T$ steps? Limits: grid up to
$1000 \times 1000$, $T \le 10^{18}$.
\solution The state is the triple (row, column, direction), so there are at most $4nm = 4\cdot10^6$ states. Store the first
time each state occurs in an array indexed by \cc{(r * m + c) * 4 + d}. At the first repetition, skip all complete cycles
that fit before $T$, then simulate the rest. Before the repetition at most $4nm$ steps are simulated, and after it fewer than
$L \le 4nm$, so the total work is $\Oh{nm}$ whatever the value of $T$.
\end{problem}

\begin{figure}[htbp]
\centering
\begin{tikzpicture}
  % left: the grid and the run (computed by simulation: tail of 5 states, cycle of 14)
  \begin{scope}[x=0.88cm, y=0.88cm]
    \foreach \rr in {0,...,3} \foreach \kk in {0,...,4}
      \draw[line width=0.4pt, draw=fgline!80, fill=white] (\kk, -\rr) rectangle ++(1, -1);
    \foreach \rr/\kk in {0/0, 0/2, 1/0}
      \draw[line width=0.4pt, draw=fgline!80, fill=fgline!55] (\kk, -\rr) rectangle ++(1, -1);
    \foreach \kk in {0,...,4} \node[lbl] at (\kk+0.5, 0.3) {\kk};
    \foreach \rr in {0,...,3} \node[lbl] at (-0.3, -\rr-0.5) {\rr};
    % the cycle through rows 2 and 3
    \draw[line width=1.6pt, draw=fgblue, rounded corners=5pt] (0.5, -2.5) -- (4.5, -2.5) -- (4.5, -3.5) -- (0.5, -3.5)
      -- cycle;
    \draw[line width=1.6pt, draw=fgblue, -{Stealth[length=2.6mm, width=2.2mm]}] (2.3, -2.5) -- (2.75, -2.5);
    \draw[line width=1.6pt, draw=fgblue, -{Stealth[length=2.6mm, width=2.2mm]}] (2.7, -3.5) -- (2.25, -3.5);
    % the tail: up to the top edge, two right turns, back down into the cycle
    \draw[line width=1.6pt, draw=fgorange, dash pattern=on 4pt off 2pt] (4.35, -1.5) -- (4.35, -0.62) arc (180:0:0.15)
      -- (4.65, -2.3);
    \draw[line width=1.6pt, draw=fgorange, -{Stealth[length=2.6mm, width=2.2mm]}] (4.35, -1.3) -- (4.35, -0.95);
    \draw[line width=1.6pt, draw=fgorange, -{Stealth[length=2.6mm, width=2.2mm]}] (4.65, -1.7) -- (4.65, -2.05);
    \fill[fgverm] (4.35, -1.5) circle (3pt);
  \end{scope}
  % right: the same run as a sequence of states
  \begin{scope}[xshift=7.1cm, yshift=-1.45cm]
    \draw[line width=1.1pt, draw=fgblue] (0, 1.0) circle (1.0);
    \foreach \k in {0,...,13} \fill[fgblue] ($(0, 1.0) + ({270 + \k*360/14}:1.0)$) circle (2.3pt);
    \draw[line width=1.1pt, draw=fgorange] (0, -2.0) -- (0, 0);
    \foreach \k in {1,...,5} \fill[fgorange] (0, {-0.4*\k}) circle (2.3pt);
    \draw[->, line width=0.7pt, draw=fgblue] (0, 1.0) ++(212:0.55) arc (212:328:0.55);
    \node[font=\sffamily\small, text=fgblue] at (0, 1.05) {$L = 14$};
    \node[lbl, anchor=west] at (0.2, -2.0) {$t = 0$ (start)};
    \node[lbl, anchor=west] at (0.2, -0.3) {$t = 5,\ 19,\ 33, \ldots$};
    \node[lbl, anchor=east] at (-0.2, -1.0) {tail};
  \end{scope}
\end{tikzpicture}
\caption{A robot run on a $4 \times 5$ grid with three walls (grey), starting at the dot facing up. It goes up, turns right
twice at the edge, comes back down (orange tail) and then circles rows 2 and 3 forever (blue cycle). Right: the same run as a
sequence of states. The state at time 5 recurs at time 19, so for every $T \ge 5$ the state at time $T$ is the state at time
$5 + (T - 5) \bmod 14$.}
\label{fig:robot}
\end{figure}

\begin{code}
void solve() {
    int n, m;
    ll T;
    cin >> n >> m >> T;
    vector<string> g(n);
    cin >> g;
    const int dr[4] = {-1, 0, 1, 0}, dc[4] = {0, 1, 0, -1};  // U, R, D, L
    int r = 0, c = 0, d = 0;
    rep(i, n) rep(j, m) if (g[i][j] == 'R') r = i, c = j;
    vector<ll> seen(4 * n * m, -1);  // first time each state (r, c, d) occurred
    ll t = 0;
    bool skipped = false;
    while (t < T) {
        int id = (r * m + c) * 4 + d;
        if (!skipped && seen[id] >= 0) {  // repeated state: skip whole cycles
            ll len = t - seen[id];
            t += (T - t) / len * len;
            skipped = true;
            if (t == T) break;
        }
        seen[id] = t;
        int nr = r + dr[d], nc = c + dc[d];
        bool ok = nr >= 0 && nr < n && nc >= 0 && nc < m && g[nr][nc] != '#';
        if (ok) r = nr, c = nc;
        else d = (d + 1) % 4;  // blocked: turn right
        t++;
    }
    print(r + 1, c + 1);
}
\end{code}

After the jump, the flag \cc{skipped} stops the program from detecting the same cycle again, and the loop simulates the
remaining steps normally. On a $1000 \times 1000$ grid with $T = 10^{18}$ the program runs in 0.08\,s. When the states are
too many to store, or are not small integers, Floyd's and Brent's cycle-finding algorithms find the same tail and cycle
lengths with $\Oh{1}$ memory by running two copies of the process at different speeds (Chapter~\ref{ch:graphs}); when many
queries ask for the state after $T$ steps from \emph{different} starting states, binary lifting answers each in
$\Oh{\log T}$ (Chapter~\ref{ch:trees}).

\begin{pitfall}[The state must be complete]
Cycle skipping is correct only if the state captures \emph{everything} that influences the future. Forgetting the direction,
a counter, a timer that changes the rules every $k$ steps, or ``whose turn it is'' makes two different situations look
identical and produces a wrong cycle, and often the samples still pass. When a step also depends on $t$ itself (for example
``every tenth step is special''), add $t \bmod 10$ to the state.
\end{pitfall}

% ======================================================================================================================
\section[Complete search with loops]{Complete search with loops \hfill \bronze}\label{sec:search-loops}

A \defn{complete search} enumerates every candidate answer, checks each one, and keeps the best valid one. Its correctness is
obvious, which is its great strength: the only questions are whether the candidates really cover every possible answer, and
whether there are few enough of them. When a candidate is described by a few numbers, nested loops are all you need: all
pairs of points, all triples of indices, all rectangles given by two opposite corners, all starting times of a schedule. The
limits tell you how many nested loops you can afford: two loops over $n = 5000$ or three over $n = 400$ are about
$2.5\cdot10^7$ and $6.4\cdot10^7$ steps.

The skill that makes complete search powerful is reducing the number of candidates \emph{without losing the optimum}. Very
often an infinite or huge space of candidates contains a small set of \emph{critical values} that is guaranteed to contain an
optimal answer, and only those need to be tried.

\begin{problem}{The most covered point}
Given $n \le 2000$ intervals $[a_i, b_i]$ with real endpoints, find a point covered by as many intervals as possible.
\solution The point can be any real number, but some optimal point is always a left endpoint $a_i$: take any optimal point
$x$ and slide it to the left until it reaches the nearest left endpoint $a_i \le x$. While sliding, the point never leaves an
interval (it would have to pass that interval's left endpoint first), so the number of covering intervals cannot decrease.
Therefore it suffices to try the $n$ left endpoints and count the intervals covering each, which is $\Oh{n^2}$
(Figure~\ref{fig:critical}). The sweep of Section~\ref{sec:events} solves the same problem in $\Oh{n \log n}$, and the
argument above is exactly why the sweep only needs to look at endpoints.
\end{problem}

\begin{figure}[htbp]
\centering
\begin{tikzpicture}[x=0.95cm, y=0.5cm]
  \numline{0}{10}{0}{}
  \ivl{1}{5}{1}{fgblue!55}{}
  \ivl{2}{6}{2}{fgblue!55}{}
  \ivl{3}{4}{3}{fgblue!55}{}
  \ivl{4}{8}{4}{fgblue!55}{}
  \ivl{7}{10}{5}{fgblue!55}{}
  \draw[line width=1pt, draw=fgverm, dashed] (4, 0.15) -- (4, 5.8) node[above, font=\sffamily\scriptsize\bfseries,
    text=fgverm] {best: 4 intervals};
  \foreach \x/\cnt in {1/1, 2/2, 3/3, 4/4, 7/2} {
    \fill[fgorange] (\x, 0) circle (2.6pt);
    \node[font=\sffamily\scriptsize\bfseries, text=fgorange!80!black] at (\x, -1.75) {\cnt};
  }
  \node[lbl, anchor=east] at (0.2, -1.75) {covered by};
\end{tikzpicture}
\caption{Five intervals. Only the left endpoints (orange dots) are candidate answers: sliding any point to the left until it
reaches a left endpoint never loses an interval. The counts below the axis are the numbers of intervals covering each
candidate; the best is 4, at $x = 4$.}
\label{fig:critical}
\end{figure}

\paragraph{How big is the search space?} Before writing a search, count its candidates and compare with the time limit.
Table~\ref{tab:searchspaces} lists the spaces that come up again and again, with the largest $n$ for which they hold about
$10^8$ candidates. Two rules of thumb follow from it: enumerating subsets is fine up to $n \approx 20$ (where each subset
costs $\Oh{n}$ to process), and enumerating permutations up to $n \approx 10$.

\begin{table}[htbp]
\centering
\small
\begin{tabular}{@{}l l c P{5.4cm}@{}}
  \toprule
  \thd{Search space} & \thd{Size} & \thd{$n$ at $\approx 10^8$} & \thd{Typical use} \\
  \midrule
  subsets of $n$ items & $2^n$ & 26 & choose any group (Apple Division) \\
  subsets of size $n/2$ & $\binom{n}{n/2}$ & 29 & split into two equal halves \\
  permutations & $n!$ & 11 & orders, schedules (Creating Strings) \\
  pairs (mask, submask) & $3^n$ & 17 & DP over subsets of subsets \\
  partitions into groups & Bell number $B_n$ & 13 & split items into any number of teams \\
  balanced bracket sequences & Catalan $C_n$ & 17 & expressions, binary tree shapes \\
  partitions of the integer $n$ & $p(n)$ & 95 & write $n$ as a sum of positive parts \\
  one half, meet in the middle & $2 \cdot 2^{n/2}$ & 54 & subset sums (Section~\ref{sec:mitm}) \\
  \bottomrule
\end{tabular}
\caption{Sizes of common search spaces. The third column is the largest $n$ for which the space has about $10^8$ elements;
if each candidate costs more than a few operations, reduce it accordingly.}
\label{tab:searchspaces}
\end{table}

\begin{trick}[Search the smaller space]
Before searching, compare the space of \emph{inputs} with the space of \emph{answers}, and search whichever is smaller.
If the answer is a small integer (a number of bits, a number of distinct values, a length up to 60), loop over the answer and
check each value. If an optimum is always attained at an input value (a median, an endpoint, a coordinate of some point),
try only those. And if a rectangle or a line is optimal, it can usually be moved until it touches input points, so only the
rectangles or lines determined by one or two input points need to be tried.
\end{trick}

% ======================================================================================================================
\section[Generating subsets]{Generating subsets \hfill \bronze}\label{sec:subsets}

A set of $n$ items has $2^n$ subsets. Represent a subset by an $n$-bit number, a \defn{bitmask}, whose bit $i$ is 1 exactly
when item $i$ belongs to the subset. Counting from $0$ to $2^n - 1$ then visits every subset exactly once
(Figure~\ref{fig:bitmasks}, left), and set operations become single machine instructions: union is \cc{a | b},
intersection \cc{a & b}, difference \cc{a & ~b}, ``is item $i$ in $S$?'' is \cc{S >> i & 1}, and the size of $S$ is
\cc{__builtin_popcount(S)}. Chapter~\ref{ch:bits} collects many more bit operations.

\begin{figure}[htbp]
\centering
\begin{tikzpicture}[b/.style={cellc, minimum width=6mm, minimum height=5.4mm, font=\ttfamily\small}]
  \node[title] at (1.9, 1.15) {counting order};
  \foreach \h/\x in {mask/0, c/1.2, b/1.8, a/2.4, subset/3.65} \node[hdr] at (\x, 0.55) {\h};
  \foreach \m/\bc/\bb/\ba/\s in {0/0/0/0/{$\emptyset$}, 1/0/0/1/{$\{a\}$}, 2/0/1/0/{$\{b\}$}, 3/0/1/1/{$\{a, b\}$},
      4/1/0/0/{$\{c\}$}, 5/1/0/1/{$\{a, c\}$}, 6/1/1/0/{$\{b, c\}$}, 7/1/1/1/{$\{a, b, c\}$}} {
    \pgfmathsetmacro{\yy}{-0.56*\m}
    \node[font=\sffamily\small\bfseries, text=fgblue] at (0, \yy) {\m};
    \edef\sA{\ifnum\ba=1 cool\else n\fi}\edef\sB{\ifnum\bb=1 cool\else n\fi}\edef\sC{\ifnum\bc=1 cool\else n\fi}
    \node[b, \sC] at (1.2, \yy) {\bc};
    \node[b, \sB] at (1.8, \yy) {\bb};
    \node[b, \sA] at (2.4, \yy) {\ba};
    \node[font=\small, anchor=west] at (3.05, \yy) {\s};
  }
  \begin{scope}[xshift=6.6cm]
    \node[title] at (2.2, 1.15) {Gray code order};
    \foreach \h/\x in {$k$/0, c/1.2, b/1.8, a/2.4, change/3.35, subset/4.45} \node[hdr] at (\x, 0.55) {\h};
    \foreach \m/\bc/\bb/\ba/\f/\ch/\s in {0/0/0/0/x/{}/{$\emptyset$}, 1/0/0/1/a/{$+a$}/{$\{a\}$},
        2/0/1/1/b/{$+b$}/{$\{a, b\}$}, 3/0/1/0/a/{$-a$}/{$\{b\}$}, 4/1/1/0/c/{$+c$}/{$\{b, c\}$},
        5/1/1/1/a/{$+a$}/{$\{a, b, c\}$}, 6/1/0/1/b/{$-b$}/{$\{a, c\}$}, 7/1/0/0/a/{$-a$}/{$\{c\}$}} {
      \pgfmathsetmacro{\yy}{-0.56*\m}
      \node[font=\sffamily\small\bfseries, text=fgblue] at (0, \yy) {\m};
      \edef\sA{\if a\f hot\else\ifnum\ba=1 cool\else n\fi\fi}
      \edef\sB{\if b\f hot\else\ifnum\bb=1 cool\else n\fi\fi}
      \edef\sC{\if c\f hot\else\ifnum\bc=1 cool\else n\fi\fi}
      \node[b, \sC] at (1.2, \yy) {\bc};
      \node[b, \sB] at (1.8, \yy) {\bb};
      \node[b, \sA] at (2.4, \yy) {\ba};
      \node[font=\small, text=fgverm] at (3.35, \yy) {\ch};
      \node[font=\small, anchor=west] at (3.9, \yy) {\s};
    }
  \end{scope}
\end{tikzpicture}
\caption{The 8 subsets of $\{a, b, c\}$ (bit 0 is $a$, bit 1 is $b$, bit 2 is $c$; blue bits are chosen items). Left: in
counting order, one step can change many bits ($3 \to 4$ changes all three). Right: in Gray code order $g_k = k \oplus
\lfloor k/2 \rfloor$, every step flips exactly one bit (red), and step $k$ flips bit number $\operatorname{ctz}(k)$.}
\label{fig:bitmasks}
\end{figure}

\begin{problem}{Apple Division (CSES)}
Divide $n \le 20$ apples with given weights into two groups so that the difference of the group weights is as small as
possible.
\solution Try all $2^n$ splits: the mask lists the apples that go into the first group, and the second group weighs the total
minus the first. With $n = 20$ that is about $10^6$ masks times 20 bits each, $2\cdot10^7$ simple operations.
\end{problem}

\begin{code}
void solve() {
    int n;
    cin >> n;
    vi p(n);
    cin >> p;
    ll total = accumulate(all(p), 0LL), best = LLONG_MAX;
    rep(mask, 1 << n) {  // the apples in the first group
        ll s = 0;
        rep(i, n) if (mask >> i & 1) s += p[i];
        chmin(best, abs(total - 2 * s));  // |group 1 - group 2|
    }
    print(best);
}
\end{code}

\usetemplate{All subsets (bitmask loop)}{Ctrl+Alt+T, 8}{cpsubsets}

\subsection{Faster subset loops}

The loop above spends $\Oh{n}$ on each subset to recompute its sum, $\Oh{n 2^n}$ in total. Two classic tricks remove the
factor $n$, and both matter when $n$ is 25 or when the work per subset is heavier than a sum.

\emph{Incremental tables.} Every non-empty mask is a smaller mask plus one item, its lowest set bit. So
\cc{sum[mask] = sum[mask & (mask - 1)] + a[__builtin_ctz(mask)]} fills a table of all $2^n$ subset sums in $\Oh{2^n}$, in
increasing order of masks, because \cc{mask & (mask - 1)} (the mask without its lowest bit) is always smaller. The same works
for any quantity that can be extended by one item: the number of edges inside a vertex set, a bitwise OR, a product modulo a
prime.

\emph{Gray code order.} When the per-subset state is a whole structure that cannot be stored $2^n$ times (a multiset, a DSU,
a board position), visit the subsets in \defn{Gray code} order instead: consecutive subsets differ in exactly one item
(Figure~\ref{fig:bitmasks}, right), so the structure is updated by a single insertion or deletion per step. The $k$-th subset
is $g_k = k \oplus \lfloor k/2\rfloor$, and going from $g_{k-1}$ to $g_k$ flips bit $\operatorname{ctz}(k)$, the number of
trailing zeros of $k$. The template generates this order.

\usetemplate{Gray code order}{Ctrl+Alt+U, 0}{cpgray}

\subsection{Submasks and $k$-subsets}

To iterate over all \defn[submask]{submasks} of a mask $S$ (the subsets of the set $S$), including $S$ itself and the empty
set, use
\cc|for (int T = S; ; T = (T - 1) & S) { ...; if (T == 0) break; }|.
Subtracting one clears the lowest set bit of $T$ and sets all the bits below it; masking with $S$ keeps only the bits that
belong to $S$, and the result is exactly the next smaller submask. The loop visits the $2^{|S|}$ submasks in decreasing
order. Summed over all masks $S$ of $n$ bits it visits $\sum_S 2^{|S|} = 3^n$ pairs, because each of the $n$ items is
independently outside $S$, in $S$ but not in $T$, or in both: this is the $3^n$ of Table~\ref{tab:searchspaces}, and the
basis of the subset DPs of Chapter~\ref{ch:dppatterns}.

\begin{trick}[Gosper's hack: all subsets of size $k$]
To visit only the $\binom{n}{k}$ masks with exactly $k$ bits (for $k \ge 1$), in increasing order, jump from each such mask
directly to the next one:
\begin{code}
for (int x = (1 << k) - 1; x < (1 << n);) {
    // ... use x ...
    int low = x & -x, r = x + low;  // r: the lowest block of ones carried
    x = (((r ^ x) >> 2) / low) | r;  // move the rest of the block to the bottom
}
\end{code}
Adding the lowest set bit \cc{low} turns the lowest block of consecutive ones into a single one bit, one position higher;
the other ones of that block are then shifted back down to the lowest positions. This is \defn{Gosper's hack}, and it beats
filtering all $2^n$ masks by popcount whenever $\binom{n}{k}$ is much smaller than $2^n$.
\end{trick}

% ======================================================================================================================
\section[Generating permutations]{Generating permutations \hfill \bronze}\label{sec:permutations}

An ordering of $n$ items is a permutation; there are $n!$ of them ($8! = 40\,320$, $10! \approx 3.6 \cdot 10^6$, $12!
\approx 4.8 \cdot 10^8$), so enumerating them is feasible up to about $n = 10$. The standard library does the work:
\cc{next_permutation} rearranges a range into the next permutation in lexicographic order and returns \cc{false} after the
last one (when it wraps around to the sorted order). Starting from the sorted order therefore visits every permutation once.
Because it compares values, a sequence with repeated values yields each \emph{distinct} arrangement exactly once: the four
letters of \texttt{aabb} give $4!/(2!\,2!) = 6$ arrangements, not 24.

The algorithm behind \cc{next_permutation} is short and worth knowing, because variants of it solve many problems (``the
next larger number with the same digits'', ``the lexicographically next valid bracket sequence''). Scan from the right for
the longest \emph{non-increasing} suffix: it is already the largest arrangement of its elements, so it cannot be increased on
its own. The element just before it, the \emph{pivot}, must grow, and by as little as possible: swap it with the smallest
element of the suffix that is larger than it. The suffix is still non-increasing after the swap, so reversing it sorts it
into its smallest arrangement (Figure~\ref{fig:nextperm}). Each call costs $\Oh{n}$ in the worst case and $\Oh{1}$ on
average over all permutations.

\begin{figure}[htbp]
\centering
\begin{tikzpicture}
  \node[blbl, anchor=east] at (-0.25, 0) {current};
  \arrs{A}{0}{0}{1/n, 3/hot, 5/cool, 4/cool, 2/cool}
  \spanabove{A-3}{A-5}{longest non-increasing suffix}
  \ptrb{A-2}{pivot}{fgverm}
  \node[notel, anchor=west, text width=6.9cm] at (4.1, 0) {The suffix $5\,4\,2$ is already the largest arrangement of
    its elements, so the pivot 3 in front of it must grow.};
  \node[blbl, anchor=east] at (-0.25, -2.2) {1.\ swap};
  \arrs{B}{0}{-2.2}{1/n, 4/hot, 5/cool, 3/cool, 2/cool}
  \draw[<->, line width=0.9pt, draw=fgverm] ($(B-2.north)+(0,0.06)$) to[out=65, in=115] ($(B-4.north)+(0,0.06)$);
  \node[notel, anchor=west, text width=6.9cm] at (4.1, -2.2) {Swap the pivot with the smallest suffix element that is
    larger than it (4). The suffix $5\,3\,2$ is still non-increasing.};
  \node[blbl, anchor=east] at (-0.25, -3.8) {2.\ reverse};
  \arrs{C}{0}{-3.8}{1/n, 4/hot, 2/good, 3/good, 5/good}
  \spanbelow{C-3}{C-5}{sorted: smallest suffix}
  \node[notel, anchor=west, text width=6.9cm] at (4.1, -3.8) {Reversing a non-increasing suffix sorts it. Result:
    $1\,3\,5\,4\,2 \to 1\,4\,2\,3\,5$, the next permutation.};
\end{tikzpicture}
\caption{How \texttt{next\_permutation} finds the next arrangement in lexicographic order: find the pivot, swap it with its
successor in the suffix, reverse the suffix. Only the pivot and the suffix change.}
\label{fig:nextperm}
\end{figure}

\begin{problem}{Creating Strings (CSES)}
Print all distinct rearrangements of a string of at most 8 letters, in alphabetical order, preceded by their number.
\end{problem}

\begin{code}
void solve() {
    string s;
    cin >> s;
    sort(all(s));
    vector<string> out;
    do out.push_back(s);
    while (next_permutation(all(s)));  // each distinct arrangement once
    print(sz(out));
    for (auto& t : out) print(t);
}
\end{code}

\usetemplate{All permutations}{Ctrl+Alt+T, 9}{cpperm}

\begin{pitfall}[Start sorted]
\cc{next_permutation} only visits the permutations that come \emph{after} the starting sequence. Forgetting to sort first
silently skips part of the search space, and the samples, which are often already sorted, may still pass. When permuting
indices rather than values, start from \cc{iota(all(p), 0)}.
\end{pitfall}

\begin{trick}[Jump straight to the $k$-th permutation]
The permutations of $0, \ldots, n-1$ in lexicographic order come in $n$ blocks of $(n-1)!$, one block per first element;
inside a block the remaining elements again form blocks of $(n-2)!$, and so on. So the $k$-th permutation (counting from 0)
takes as its first element the $\lfloor k/(n-1)! \rfloor$-th smallest unused value, and continues with $k \bmod (n-1)!$.
The successive quotients are the digits of $k$ in the \defn{factorial number system} (the \defn{Lehmer code} of the
permutation), and running the procedure backwards gives the rank of a permutation. Both take $\Oh{n^2}$ as written, or
$\Oh{n \log n}$ with a Fenwick tree that finds the $d$-th unused value (Chapter~\ref{ch:segtree}).
\begin{code}
vi kth_perm(int n, ll k) {  // n <= 20
    vector<ll> f(n + 1, 1);
    rep(i, 1, n + 1) f[i] = f[i - 1] * i;
    vi pool(n), out;
    iota(all(pool), 0);
    for (int i = n - 1; i >= 0; i--) {
        int d = k / f[i];  // skip d whole blocks of i! permutations
        k %= f[i];
        out.push_back(pool[d]);
        pool.erase(pool.begin() + d);
    }
    return out;
}
\end{code}
\end{trick}

% ======================================================================================================================
\section[Backtracking and pruning]{Backtracking and pruning \hfill \silver}\label{sec:backtracking}

\defn{Backtracking} builds a solution one choice at a time and abandons a partial solution as soon as it is clear that it
cannot be completed into a valid (or better) one. The partial solutions form a \emph{search tree}: the root is the empty
solution, the children of a node are its extensions by one more choice, and the leaves are complete candidates. A plain loop
over all candidates visits every leaf; backtracking walks the tree depth first and \emph{prunes}, that is, refuses to enter a
subtree once it knows that the subtree contains nothing useful. A single pruned node near the root removes an exponential
number of leaves, which is why the difference between a search that runs in a second and one that runs for a year is almost
always \defn{pruning}: how early, and how often, the search can say no.

\begin{pseudo}{Backtracking}
\begin{algorithmic}[1]
\Procedure{Search}{$\mathit{depth}$}
  \If{the partial solution is complete}
    \State record it (count it, or compare it with the best so far)
    \State \Return
  \EndIf
  \ForAll{choices $x$ for the next position}
    \If{$x$ is consistent with the choices so far and can still win}
      \State apply $x$ to the shared state \Comment{mark a column, add a weight, \ldots}
      \State \Call{Search}{$\mathit{depth} + 1$}
      \State undo $x$ \Comment{restore the state exactly}
    \EndIf
  \EndFor
\EndProcedure
\end{algorithmic}
\end{pseudo}

The \emph{undo} line is the heart of backtracking. The search keeps one shared state, modifies it before descending and
restores it after returning, so no copies are ever made and each step costs $\Oh{1}$. Forgetting to undo, or undoing only part
of the change, is the classic backtracking bug: the state after returning must be \emph{identical} to the state before
descending.

\begin{problem}{Chessboard and Queens (CSES)}
Place eight queens on an $8 \times 8$ board, avoiding reserved squares, so that no two queens attack each other. Count the
ways.
\solution Two queens in the same row attack each other, so every row holds exactly one queen: choose the column of the
queen in row 0, then in row 1, and so on. A queen in row $r$, column $c$ attacks its column, its diagonal (on which $r + c$ is
constant) and its anti-diagonal (on which $r - c$ is constant). Three boolean arrays record which columns and diagonals are
taken, so each square is checked in $\Oh{1}$. Figure~\ref{fig:queens} shows the complete search tree for four queens: the
search visits 16 placements, while the tree of all ways to put one queen in each row has $4 + 16 + 64 + 256 = 340$ nodes.
\end{problem}

\begin{figure}[htbp]
\centering
\begin{minipage}[c]{0.55\linewidth}
\centering
\begin{forest}
  for tree={circle, draw=fgink!80, fill=fgfill, line width=0.7pt, minimum size=5.6mm, inner sep=0pt,
    font=\sffamily\scriptsize, s sep=2.4mm, l sep=4.5mm, edge={line width=0.7pt, draw=fgline!80!fgink}}
  [{}, fill=fgline!40
    [1 [3, hot] [4 [2, hot]]]
    [2 [4 [1 [3, good]]]]
    [3 [1 [4 [2, good]]]]
    [4 [1 [3, hot]] [2, hot]]]
\end{forest}
\end{minipage}%
\begin{minipage}[c]{0.3\linewidth}
\centering
\begin{tikzpicture}[scale=0.62]
  \foreach \rr in {0,...,3} \foreach \kk in {0,...,3} {
    \pgfmathtruncatemacro{\cpar}{mod(\rr+\kk,2)}
    \ifnum\cpar=0 \fill[fgline!30] (\kk, -\rr) rectangle ++(1, -1); \else \fill[white] (\kk, -\rr) rectangle ++(1, -1); \fi
  }
  \draw[line width=0.6pt, draw=fgink!70] (0, 0) rectangle (4, -4);
  \foreach \rr/\kk in {0/1, 1/3, 2/0, 3/2}
    \node[circle, fill=fggreen, minimum size=4.6mm, inner sep=0pt, text=white, font=\sffamily\tiny\bfseries]
      at (\kk+0.5, -\rr-0.5) {Q};
  \foreach \kk in {1,...,4} \node[lbl] at (\kk-0.5, 0.35) {\kk};
  \node[lbl] at (2, -4.55) {columns $2, 4, 1, 3$};
\end{tikzpicture}
\end{minipage}
\caption{Left: the whole backtracking tree for four queens; a node is the column of the queen placed in the next row. Red
nodes are dead ends (the following row has no safe square), green nodes complete a solution. Right: the solution
$2, 4, 1, 3$; the other one is its mirror image $3, 1, 4, 2$.}
\label{fig:queens}
\end{figure}

\begin{code}
void solve() {
    vector<string> b(8);
    cin >> b;
    array<bool, 8> col{};
    array<bool, 15> diag{}, anti{};  // r + c, and r - c + 7
    ll ways = 0;
    auto place = [&](auto&& self, int r) -> void {  // one queen per row
        if (r == 8) {
            ways++;
            return;
        }
        rep(c, 8) {
            bool taken = col[c] || diag[r + c] || anti[r - c + 7];
            if (b[r][c] == '*' || taken) continue;
            col[c] = diag[r + c] = anti[r - c + 7] = true;
            self(self, r + 1);
            col[c] = diag[r + c] = anti[r - c + 7] = false;  // undo
        }
    };
    place(place, 0);
    print(ways);
}
\end{code}

\usetemplate{Recursive complete search}{Ctrl+Alt+T, Shift+Y}{cpbacktrack}

\subsection{A case study in pruning: Grid Paths}

\begin{problem}{Grid Path Description (CSES)}
A path on a $7 \times 7$ grid starts in the upper-left corner, ends in the lower-left corner, and visits every one of the 49
squares exactly once, so it consists of 48 moves. You are given a description of 48 characters, each \texttt{U}, \texttt{D},
\texttt{L}, \texttt{R} (the move is fixed) or \texttt{?} (any move). Count the paths that match the description.
\end{problem}

A naive search tries up to three directions at each of 48 steps, $3^{47} \approx 2.7 \cdot 10^{22}$ sequences in the worst
case (a description of 48 question marks). Refusing to step onto visited squares already cuts this enormously, but the
search still visits $1.9 \cdot 10^{10}$ nodes and needs two minutes (Table~\ref{tab:gridpruning}). Two observations about
\emph{dead} partial paths bring it down to a quarter of a second (Figure~\ref{fig:gridpaths}).

\begin{enumerate}
  \item \textbf{The end is reached too early.} Every path must end at the lower-left corner, so a path that arrives there
        before its 48th move can never be completed: stop at once instead of wandering on.
  \item \textbf{The free squares split in two.} Suppose the squares above and below the current square are blocked (visited
        or outside the grid) while the squares to its left and right are both free, or the same with the two directions
        exchanged. The blocked squares, the path and the border then form a barrier through the current square, and the
        free squares on its two sides are disconnected from each other. The path can continue into only one side and can
        never come back, so the squares on the other side will never be visited: the partial path is dead.
\end{enumerate}

The second rule fires far more often than one would guess: every time the path runs straight into a wall or into its own
earlier trail with space on both sides. Because it cuts subtrees close to the root, it removes 99.6\% of the remaining search.

\begin{figure}[htbp]
\centering
\begin{tikzpicture}[x=0.46cm, y=0.46cm]
  \tikzset{gpath/.style={line width=1.6pt, draw=fgverm}, gvis/.style={fill=fgline!35},
    gend/.style={draw=fggreen!80!black, line width=1.1pt}, gcur/.style={draw=fgverm, line width=1.1pt}}
  % (a) the end corner is reached after 6 moves
  \begin{scope}
    \fill[fgsky!30] (0, 0) rectangle (7, -7);
    \foreach \rr in {0,...,6} \fill[gvis] (0, -\rr) rectangle ++(1, -1);
    \draw[fgline!80, line width=0.4pt, step=1] (0, 0) grid (7, -7);
    \draw[gpath] (0.5, -0.5) -- (0.5, -6.5);
    \fill[fgverm] (0.5, -0.5) circle (2.6pt);
    \draw[gend] (0.5, -6.5) circle (6.5pt);
    \node[font=\sffamily\scriptsize\bfseries, text=fgblue] at (4, -3.5) {42 unvisited};
    \node[lbl, text width=3.6cm, align=center, anchor=north] at (3.5, -7.4) {(a) the end corner is reached after
      6 moves: stop};
  \end{scope}
  % (b) the path runs into the bottom wall with free squares on both sides
  \begin{scope}[xshift=4.3cm]
    \fill[fgsky!30] (0, -1) rectangle (3, -7);
    \fill[fgorange!35] (4, 0) rectangle (7, -7);
    \foreach \kk in {0,...,3} \fill[gvis] (\kk, 0) rectangle ++(1, -1);
    \foreach \rr in {1,...,6} \fill[gvis] (3, -\rr) rectangle ++(1, -1);
    \draw[fgline!80, line width=0.4pt, step=1] (0, 0) grid (7, -7);
    \draw[gpath] (0.5, -0.5) -- (3.5, -0.5) -- (3.5, -6.5);
    \fill[fgverm] (0.5, -0.5) circle (2.6pt);
    \draw[gend] (0.5, -6.5) circle (6.5pt);
    \draw[gcur] (3, -6) rectangle (4, -7);
    \node[font=\sffamily\scriptsize\bfseries, text=fgblue] at (1.5, -3.5) {18};
    \node[font=\sffamily\scriptsize\bfseries, text=fgorange!80!black] at (5.5, -3.5) {21};
    \node[lbl, text width=3.6cm, align=center, anchor=north] at (3.5, -7.4) {(b) blocked below, free left and right:
      one side is lost};
  \end{scope}
  % (c) the path encloses a pocket against its own trail
  \begin{scope}[xshift=8.6cm]
    \fill[fgsky!30] (0, 0) rectangle (7, -7);
    \foreach \rr/\kk in {0/0, 1/0, 2/0, 3/0, 3/1, 3/2, 3/3, 3/4, 3/5, 2/5, 1/5, 1/4, 1/3, 2/3}
      \fill[gvis] (\kk, -\rr) rectangle ++(1, -1);
    \fill[fgorange!35] (4, -2) rectangle (5, -3);
    \draw[fgline!80, line width=0.4pt, step=1] (0, 0) grid (7, -7);
    \draw[gpath] (0.5, -0.5) -- (0.5, -3.5) -- (5.5, -3.5) -- (5.5, -1.5) -- (3.5, -1.5) -- (3.5, -2.5);
    \fill[fgverm] (0.5, -0.5) circle (2.6pt);
    \draw[gend] (0.5, -6.5) circle (6.5pt);
    \draw[gcur] (3, -2) rectangle (4, -3);
    \node[lbl, text width=3.6cm, align=center, anchor=north] at (3.5, -7.4) {(c) blocked above and below by its own
      trail: the pocket is lost};
  \end{scope}
\end{tikzpicture}
\caption{Dead partial paths in Grid Paths. The path (red) starts at the dot; the green circle marks the end corner and the red
square the current position. Blue squares are free; orange squares are free but can no longer be reached once the path
chooses a side. In (b) and (c) the current square is blocked on two opposite sides and free on the other two, which is
exactly the test in the code.}
\label{fig:gridpaths}
\end{figure}

\begin{table}[htbp]
\centering
\small
\begin{tabular}{@{}l r r@{}}
  \toprule
  \thd{Pruning rules} & \thd{Search nodes} & \thd{Time} \\
  \midrule
  never step onto a visited square & $1.9 \cdot 10^{10}$ & 120\,s \\
  $+$ stop when the end corner is reached & $1.1 \cdot 10^{10}$ & 54\,s \\
  $+$ stop when the free squares split in two & $4.7 \cdot 10^{7}$ & 0.24\,s \\
  \bottomrule
\end{tabular}
\caption{Counting all 88\,418 paths (a description of 48 question marks) with more and more pruning, measured with
\texttt{g++ -O2}. The split rule alone makes the search 230 times faster.}
\label{tab:gridpruning}
\end{table}

\begin{code}
void solve() {
    string s;
    cin >> s;
    // the 7 x 7 grid is rows and columns 1..7, surrounded by blocked cells
    array<array<bool, 9>, 9> vis{};
    rep(i, 9) vis[i][0] = vis[i][8] = vis[0][i] = vis[8][i] = true;
    const int dr[4] = {1, -1, 0, 0}, dc[4] = {0, 0, -1, 1};
    const string dirs = "DULR";
    ll ways = 0;
    auto go = [&](auto&& self, int r, int c, int step) -> void {
        if (r == 7 && c == 1) {  // every path ends here, so stop here
            if (step == 48) ways++;
            return;
        }
        // blocked on two opposite sides, free on the other two: the free
        // cells split into two parts and one of them can never be visited
        bool up = vis[r - 1][c], down = vis[r + 1][c];
        bool left = vis[r][c - 1], right = vis[r][c + 1];
        if (up && down && !left && !right) return;
        if (left && right && !up && !down) return;
        vis[r][c] = true;
        rep(k, 4) {
            if (s[step] != '?' && s[step] != dirs[k]) continue;
            int nr = r + dr[k], nc = c + dc[k];
            if (!vis[nr][nc]) self(self, nr, nc, step + 1);
        }
        vis[r][c] = false;
    };
    go(go, 1, 1, 0);
    print(ways);
}
\end{code}

The border of blocked cells makes every neighbour access safe and turns ``outside the grid'' into ``visited'', which is
exactly what both pruning rules need. The search never reads \cc{s[48]}: a path that has made 48 moves has visited all 49
squares, including the end corner, and the search stops as soon as it reaches that corner.

\begin{trick}[Pruning that matters]
\begin{itemize}
  \item \textbf{Check constraints early.} Test a choice when it is made, not when the solution is complete; test the most
        restrictive constraints first.
  \item \textbf{Look for dead states}, as in Grid Paths: parity, connectivity of what remains, or a resource that can no
        longer suffice (``the remaining items cannot fill the remaining capacity'').
  \item \textbf{Bound optimistically.} If the current cost plus a lower bound on the cost of the rest cannot beat the best
        solution so far, stop. This is \defn{branch and bound}; a cheap but tight bound is worth more than a clever search
        order.
  \item \textbf{Order the choices} so that good solutions are found first: the better the first solution, the more the
        bound prunes. Greedy order (largest item first, nearest city first) is a good default.
  \item \textbf{Choose the most constrained variable next}: in Sudoku-like problems, fill the cell with the fewest remaining
        options, so that contradictions surface near the root.
  \item \textbf{Break symmetry.} Fix the first element, or require choices in increasing order when their order does not
        matter; this divides the search by the size of the symmetry group (up to $n!$).
  \item \textbf{Memoise.} If the same partial state can be reached along different paths, you are really doing dynamic
        programming (Part~V), and caching the result of each state turns exponential time into polynomial.
\end{itemize}
\end{trick}

% ======================================================================================================================
\section[Meet in the middle]{Meet in the middle \hfill \gold}\label{sec:mitm}

For $n$ around 40, the $2^{40} \approx 10^{12}$ subsets are far too many, but $2^{20} \approx 10^6$ is nothing.
\defn{Meet in the middle} splits the items into two halves, enumerates the subsets of each half separately, and combines the
two lists with sorting and two pointers, binary search or a hash map. Every subset of all the items is a pair (subset of the
left half, subset of the right half), so questions about the whole set become questions about pairs from two lists of
$2^{n/2}$ elements each, and those can be answered in about $2^{n/2}$ time instead of $2^n$.

\begin{problem}{Meet in the Middle (CSES)}
Count the subsets of $n \le 40$ numbers (each up to $10^9$) whose sum is exactly $x$.
\solution Enumerate the subset sums $L$ of the first half and $R$ of the second half. The answer is the number of pairs
$(a, b) \in L \times R$ with $a + b = x$. With both lists sorted, walk one pointer up $L$ and another down $R$: if
$a + b < x$, the current $a$ is too small for every remaining $b$, so move up in $L$; if $a + b > x$, move down in $R$; if
they are equal, the runs of equal values on both sides contribute the product of their lengths (Figure~\ref{fig:mitm}).
\end{problem}

\begin{figure}[htbp]
\centering
\begin{tikzpicture}
  \node[title, anchor=west] at (-0.55, 1.0) {left half $3, 5$; right half $2, 7$; target $x = 10$};
  \node[blbl, anchor=east] at (-0.2, 0) {$L$};
  \arrs{L}{0}{0}{0/n, 3/hot, 5/n, 8/good}
  \node[blbl, anchor=east] at (-0.2, -1.5) {$R$};
  \arrs{R}{0}{-1.5}{0/n, 2/good, 7/hot, 9/n}
  \draw[line width=1.1pt, draw=fgverm] (L-2.south) -- (R-3.north);
  \draw[line width=1.1pt, draw=fggreen] (L-4.south) -- (R-2.north);
  \node[notel, anchor=west, text width=7cm] at (3.6, -0.75) {$3 + 7 = 10$ is the subset $\{3, 7\}$, and $8 + 2 = 10$ is
    the subset $\{3, 5, 2\}$. The pointer on $L$ only moves right and the one on $R$ only moves left, so once both lists
    are sorted the matching takes linear time.};
\end{tikzpicture}
\caption{Meet in the middle for subset sums. $L$ and $R$ hold the $2^{n/2}$ subset sums of each half, sorted. A subset of
all items is one sum from each list, so the pairs adding up to $x$ are exactly the subsets with sum $x$.}
\label{fig:mitm}
\end{figure}

\begin{code}
void solve() {
    int n;
    ll x;
    cin >> n >> x;
    vi t(n);
    cin >> t;
    auto sums = [&](int lo, int hi) {  // all subset sums of t[lo..hi), sorted
        vector<ll> s{0};
        rep(i, lo, hi) {  // merge s with s + t[i]: stays sorted, linear time
            vector<ll> shifted(sz(s)), merged(2 * sz(s));
            rep(j, sz(s)) shifted[j] = s[j] + t[i];
            merge(all(s), all(shifted), merged.begin());
            s = merged;
        }
        return s;
    };
    vector<ll> L = sums(0, n / 2), R = sums(n / 2, n);
    ll ways = 0;
    int i = 0, j = sz(R) - 1;  // two pointers: L ascending, R descending
    while (i < sz(L) && j >= 0) {
        if (L[i] + R[j] < x) i++;
        else if (L[i] + R[j] > x) j--;
        else {  // equal: multiply the sizes of the two runs of equal values
            ll a = L[i], b = R[j], ca = 0, cb = 0;
            while (i < sz(L) && L[i] == a) i++, ca++;
            while (j >= 0 && R[j] == b) j--, cb++;
            ways += ca * cb;
        }
    }
    print(ways);
}
\end{code}

\begin{trick}[Generate the sums already sorted]
Sorting $2^{20}$ numbers costs a factor of 20. Instead, build the sums one item at a time: if $s$ is the sorted list of sums
of the first $i$ items, the list for $i + 1$ items is the \emph{merge} of $s$ and $s + t_i$, two sorted lists, in linear
time. The sizes double at each step, so the whole half costs $\Oh{2^{n/2}}$, and the program above runs in 0.1\,s for
$n = 40$.
\end{trick}

\usetemplate{Meet in the middle (subset sums)}{Ctrl+Alt+T, Shift+U}{cpmitm}

\paragraph{Variations.} The combining step depends on the question. To find the subset sum closest to $x$, binary search
each $a \in L$ for $x - a$ in $R$. For subsets with a given \emph{XOR}, count matching values with a hash map. To count
subsets with sum \emph{at most} $x$, count for each $a$ the elements of $R$ that are at most $x - a$, again with two
pointers. Beyond subsets, the idea applies to any search of depth $d$ whose two halves can be combined: a search of $b^d$
states becomes two of $b^{d/2}$ each. This is how bidirectional BFS finds shortest paths in huge implicit graphs
(Chapter~\ref{ch:graphs}), and how 4-SUM ($a + b + c + d = S$) becomes pairs $(a, b)$ against pairs $(c, d)$ in
$\Oh{n^2 \log n}$.

\begin{recognize}
$n \le 40$ (or 36, or 45) together with a task that is exponential in $n$: subset sums, subset XORs, ``split into two
groups'', knapsack with huge weights. More generally, any search where $b^{d/2}$ states are affordable but $b^d$ are not, and
where a solution can be cut into two independent halves that meet in a simple condition.
\end{recognize}

% ======================================================================================================================
\section[Fracturing search]{Fracturing search: the $k$ best solutions \hfill \ioi}\label{sec:fracturing}

Sometimes you need not only the best solution but the $k$ best: the $k$ smallest subset sums, the $k$ cheapest spanning
trees, the $k$ shortest paths. \defn{Fracturing search} finds them by arranging \emph{all} solutions into a tree in which
every child is no better than its parent and every solution appears exactly once. A priority queue then produces solutions
in order: pop the best node, report it, push its children. The first pop is the root, the best solution overall; each later
pop is the best node whose parent has already been reported, and because children are never better than their parents,
that is the best solution not yet reported. After $k$ pops you have the $k$ best solutions, and the queue never holds more
than $k \times (\text{number of children}) + 1$ nodes.

\begin{problem}{The $k$ smallest subset sums}
Given $n$ positive numbers, list the $k$ smallest sums over all non-empty subsets.
\solution Sort the numbers and describe a subset by its sum and the index $i$ of its largest element. Every subset other than
$\{a_0\}$ has exactly one \emph{parent}, from which it arises by one of two moves: \emph{add} the element $a_{i+1}$, or
\emph{replace} the largest element $a_i$ by $a_{i+1}$. (To find the parent of a subset with largest index $j$: if index
$j - 1$ is also in the subset, remove $j$; otherwise replace $j$ by $j - 1$.) Neither move decreases the sum, because the
numbers are sorted and positive, so the tree rooted at $\{a_0\}$ is heap-ordered, and popping the queue $k$ times lists the
$k$ smallest sums (Figure~\ref{fig:fracturing}).
\end{problem}

\begin{figure}[htbp]
\centering
\begin{forest}
  for tree={rectangle, rounded corners=2pt, draw=fgink!70, fill=fgfill, line width=0.6pt, font=\sffamily\scriptsize,
    align=center, inner sep=2.5pt, s sep=3mm, l sep=6mm, edge={line width=0.7pt, draw=fgline!80!fgink}}
  [{$\{1\}$\\ 1}, good
    [{$\{1,2\}$\\ 3}, good, edge label={node[midway, left, font=\sffamily\tiny, text=fgblue]{add}}
      [{$\{1,2,4\}$\\ 7}, good, edge label={node[midway, left, font=\sffamily\tiny, text=fgblue]{add}}
        [{\ldots}, draw=none, fill=none]]
      [{$\{1,4\}$\\ 5}, good, edge label={node[midway, right, font=\sffamily\tiny, text=fgblue]{replace}}
        [{\ldots}, draw=none, fill=none]]]
    [{$\{2\}$\\ 2}, good, edge label={node[midway, right, font=\sffamily\tiny, text=fgblue]{replace}}
      [{$\{2,4\}$\\ 6}, good [{\ldots}, draw=none, fill=none]]
      [{$\{4\}$\\ 4}, good [{$\{4,7\}$\\ 11}] [{$\{7\}$\\ 7}, good]]]]
\end{forest}
\caption{The successor tree for $a = 1, 2, 4, 7$; each node shows a subset and its sum. Left children add the next element,
right children replace the largest element by the next one. Sums never decrease downwards, so a priority queue pops them in
order: $1, 2, 3, 4, 5, 6, 7, 7, \ldots$ (green: the first eight pops).}
\label{fig:fracturing}
\end{figure}

\begin{code}
void solve() {
    int n, k;
    cin >> n >> k;
    vi a(n);
    cin >> a;
    sort(all(a));
    minpq<pll> pq;  // (sum, i): a subset whose largest element is a[i]
    pq.push({a[0], 0});
    vi out;
    while (sz(out) < k) {
        auto [s, i] = pq.top();
        pq.pop();
        out.push_back(s);
        if (i + 1 < n) {  // two children; together they reach every subset once
            pq.push({s + a[i + 1], i + 1});         // also take a[i + 1]
            pq.push({s - a[i] + a[i + 1], i + 1});  // swap a[i] for a[i + 1]
        }
    }
    print(out);
}
\end{code}

The general recipe: find a way to \emph{fracture} the solution space into disjoint parts, each described by a few constraints
and each with an easily computed best element, such that removing the best element of a part leaves a few smaller parts of
the same form. Then let the priority queue explore the parts in order of their best elements. For the $k$ shortest paths, a
part is ``paths that share this prefix and then avoid these edges''; for the $k$ smallest spanning trees it is ``trees that
contain these edges and avoid those''. The total cost is $\Oh{k \cdot (c \cdot B + \log k)}$, where $c$ is the number of
children per node and $B$ the cost of finding the best element of a part.

% ======================================================================================================================
\section[Casework and ad hoc problems]{Casework and ad hoc problems \hfill \bronze}\label{sec:adhoc}

Many Bronze problems and most Codeforces Division~2 A and B problems have no named algorithm. The solution is a single
observation, or a careful division into cases, and such problems reward one specific habit: \emph{compute small cases with a
brute force and look at the table}. The brute force is a complete search, which you can write in minutes and trust; the table
shows you the rule; and the rule is usually much easier to prove once you know what it is.

\begin{problem}{Coin Piles (CSES)}
Two piles hold $a$ and $b$ coins. One move removes one coin from one pile and two coins from the other. Can both piles be
emptied? Answer up to $10^5$ such questions with $a, b \le 10^9$.
\solution Tabulate the answers for small $a$ and $b$ with a brute force (Figure~\ref{fig:coinpiles}). Two facts jump out, and
both are easy to prove. Every move removes three coins, so $a + b$ must be divisible by 3. And every move removes at most two
coins from one pile while removing at least one from the other, so if one pile holds more than twice the other, the larger
pile can never be emptied. The table suggests that these two conditions are also sufficient, and indeed they are: if both
hold, then $x = (2a - b)/3$ moves of the first kind and $y = (2b - a)/3$ of the second kind are non-negative integers that
empty both piles.
\end{problem}

\begin{figure}[htbp]
\centering
\begin{tikzpicture}[s/.style={minimum size=4.6mm, inner sep=0pt, draw=white, line width=0.5pt}]
  \foreach \a in {0,...,12} \foreach \b in {0,...,12} {
    \pgfmathtruncatemacro{\ok}{(mod(\a+\b,3)==0 && 2*min(\a,\b)>=max(\a,\b)) ? 1 : 0}
    \ifnum\ok=1 \node[s, fill=fggreen!55] at (\a*0.46, -\b*0.46) {};
    \else \node[s, fill=fgline!22] at (\a*0.46, -\b*0.46) {}; \fi
  }
  \foreach \a in {0, 3, 6, 9, 12} \node[lbl] at (\a*0.46, 0.42) {\a};
  \foreach \b in {0, 3, 6, 9, 12} \node[lbl] at (-0.45, -\b*0.46) {\b};
  \node[blbl] at (2.76, 0.85) {$a$};
  \node[blbl] at (-0.9, -2.76) {$b$};
  \draw[line width=0.9pt, draw=fgverm, dashed] (0, 0) -- (3.0, -6.0) node[below, font=\sffamily\scriptsize,
    text=fgverm] {$b = 2a$};
  \draw[line width=0.9pt, draw=fgverm, dashed] (0, 0) -- (6.0, -3.0) node[right, font=\sffamily\scriptsize,
    text=fgverm] {$a = 2b$};
  \node[notel, anchor=west, text width=5.4cm] at (7.5, -2.76) {Green: both piles can be emptied. The green cells lie on
    every third anti-diagonal ($a + b \equiv 0 \pmod 3$) and inside the wedge between the lines $b = 2a$ and $a = 2b$.};
\end{tikzpicture}
\caption{Coin Piles for $0 \le a, b \le 12$, computed by brute force. Looking at such a table is often the fastest way to
find the rule; proving it afterwards takes two lines.}
\label{fig:coinpiles}
\end{figure}

\begin{code}
void solve() {
    int t;
    cin >> t;
    while (t--) {
        ll a, b;
        cin >> a >> b;
        // each move removes 3 coins; no pile may exceed twice the other
        yn((a + b) % 3 == 0 && 2 * min(a, b) >= max(a, b));
    }
}
\end{code}

\begin{tip}[Getting casework right]
List the cases on paper before coding, and check that they cover every input and do not overlap. For each case, construct a
small example and check your formula on it by hand. The special values ($0$, $1$, all elements equal, $n = 1$, $n = 2$, the
maximum value) are where casework breaks, so test each of them. When the analysis grows beyond five cases, look for a simpler
characterisation, or brute force a small parameter instead of reasoning about it.
\end{tip}

\begin{trick}[Ad hoc habits]
Look for invariants (a sum, a parity or a remainder that no operation changes), extremal elements (what must happen to the
maximum?), and better representations (differences $a_{i+1} - a_i$ instead of values, prefix sums instead of elements,
positions instead of values). Above all, write the brute force, print a table of answers for all small inputs, and look at it
from every angle: by rows, by columns, by diagonals, modulo small numbers. The pattern is often the whole solution, and the
brute force then doubles as the checker for your stress test (Chapter~\ref{ch:patterns}).
\end{trick}

\begin{summary}[Chapter summary]
\begin{itemize}
  \item Simulate with an explicit state and one step function; when $T$ is huge, find the cycle and skip it. The state must
        contain everything that influences the future.
  \item Count the search space first (Table~\ref{tab:searchspaces}), search the smaller of inputs and answers, and look for
        critical values.
  \item Subsets are bitmasks ($n \le 20$); incremental tables and Gray codes remove the factor $n$; submask loops cost $3^n$
        in total; Gosper's hack visits the $k$-subsets.
  \item Permutations come from \texttt{next\_permutation} ($n \le 10$), started from the sorted order; the factorial number
        system jumps to the $k$-th one.
  \item Backtracking modifies one shared state and undoes each change. Dead-state detection, bounds, good ordering and
        symmetry breaking can make it faster by factors of hundreds.
  \item Meet in the middle turns $2^n$ into $2^{n/2}$; build the half sums already sorted by merging.
  \item Fracturing search lists the $k$ best solutions with a priority queue over a heap-ordered tree of all solutions.
  \item For ad hoc problems, tabulate small cases with a brute force and find the rule in the table.
\end{itemize}
\end{summary}

\begin{practice}
\ppgroup{Simulation and enumeration}
\begin{pplist}
  \item Number Spiral
  \item Two Knights
  \item Apple Division
  \item Creating Strings
  \item Gray Code
  \item Tower of Hanoi
\end{pplist}
\ppgroup{Backtracking, meet in the middle and ad hoc}
\begin{pplist}
  \item Chessboard and Queens
  \item Grid Path Description
  \item Meet in the Middle
  \item Coin Piles
  \item Two Sets
  \item Palindrome Reorder
\end{pplist}
\end{practice}
